\ssr{Введение}

Цели данной лабораторной работы:
\begin{enumerate}
	\item изучить основные концепции и принципы функционирования дисков и файловых систем;
	\item изучить методы управления дисками и файловыми системами в Linux;
	\item приобрести навыки работы с утилитами управления дисками и файловыми системами. 
\end{enumerate}

\ssr{Выполнение}

\ssection{Задание 1. Разметка диска, создание файловой системы} 

\begin{figure}[!htbp]
	\centering
	%\includegraphics[width=0.6\linewidth]{""}
	\caption{Разметка диска}
\end{figure}
\FloatBarrier

\ssection{Задание 2. Создание RAID-массива} 

\begin{figure}[!htbp]
	\centering
	%\includegraphics[width=0.6\linewidth]{""}
	\caption{Созданный RAID-массив}
\end{figure}
\FloatBarrier

\ssection{Задание 3. Получение информации о файловой системе} 

\begin{figure}[!htbp]
	\centering
	%\includegraphics[width=0.6\linewidth]{""}
	\caption{Тип основной файловой системы, ее размер, общее количество и количество занятых inode}
\end{figure}
\FloatBarrier

\ssection{Дополнительное задание}

\begin{figure}[!htbp]
	\centering
	%\includegraphics[width=0.6\linewidth]{""}
	\caption{Расширение файловой системы}
\end{figure}
\FloatBarrier

\ssr{Вывод}

В ходе лабораторной работы были:  
\begin{itemize}
	\item изучены основные концепции и принципы функционирования дисков и файловых систем;
	\item изучены методы управления дисками и файловыми системами в Linux;
	\item приобретены навыки работы с утилитами управления дисками и файловыми системами. 
\end{itemize}

\ssr{Контрольные вопросы}

\begin{enumerate}
	\item \textit{Дайте определение дисковой подсистемы операционной системы. Какие ее ключевые функции?}\newline 
	Совокупность аппаратных и программных компонентов, которые обеспечивают возможность доступа к физическим носителям данных и управления информацией, которая на них хранится.
	Она хранит и обрабатывает данные на физических дисках.
	\item \textit{С помощью каких команд в Linux можно получить общий список блочных устройств (дисков, разделов) и информацию о них? Какие еще типы устройств существуют? Приведите примеры}\newline 
	lsblk, fdisk -l, parted -l, blkid, ls -l /dev/disk/by-*.
	Символьные (character, /dev/tty) и сетевые (network, имеют собственные имена).
	\item \textit{В чем заключается основное назначение разделов на жестком диске?}\newline 
	%Это логически выделенные части накопителя данных, состоящие из смежных блоков, которые отображаются в системе как отдельные диски.
	Каждый раздел может иметь свою собственную файловую систему и параметры форматирования, появляется возможность установить несколько операционных систем на один жесткий диск.
	Использование разделов в сочетании с технологиями LVM и RAID
	увеличивает надежность хранения данных и снижает риск потери информации в случае сбоя в системе.
	\item \textit{Назовите ключевые различия между таблицами разделов MBR и GPT (количество разделов, максимальный размер диска, связь с BIOS/UEFI).}\newline 
	\begin{tabularx}{0.9\textwidth}{|X|X|p{4cm}|X|}
		\hline
		\textbf{Название} & \textbf{Количество разделов} & \textbf{Максимальный размер диска} & \textbf{Связь с BIOS/UEFI} \\\hline
		MBR & 4 (3+1) & 2 ТБ & BIOS\\\hline
		GPT & $\infty$ &  9.4 ЗБ & UEFI \\\hline
	\end{tabularx}
	\item \textit{Какая утилита в Linux является стандартным инструментом для работы с разделами на дисках с MBR, а какая — для GPT?}\newline 
	fdisk -- утилита для MBR, gdisk -- для GPT, parted -- для обоих.
	\item \textit{Что такое RAID? Какова основная цель использования RAID 1 (зеркалирование) и RAID 0 (страйпинг)? В чем их принципиальное различие с точки зрения надежности?}\newline 
	RAID (Redundant Array of Independent Disks) -- технология объединения нескольких физических дисков в один логический массив для повышения отказоустойчивости и/или производительности хранилища данных. 
	RAID 1 позволяет сохранить данные в случае отказа одного из дисков, при этом стоимость хранения возрастает в 2 раза по сравнению с обычным хранением данных; RAID 0 разбивает данные на несколько дисков, что улучшает скорость работы, но если один из
	дисков выйдет из строя, все данные будут утеряны.
	\item \textit{Какие виды RAID-массивов, помимо рассмотренных существуют? Что представляют собой RAID 10 и RAID 01}\newline 
	% RAID 10 = mirror of stripes (RAID 1+0): data striped across mirrored pairs. Good performance and redundancy; survives one disk per mirror; rebuild is fast.
	RAID 10 (RAID 1+0) -- зеркалирование страйпинга (данные страйпятся по парам зеркал); высокая надежность благодаря тому, что массив выдерживает отказ одного диска в каждой паре.
	RAID 01 (RAID 0+1) -- страйпинг зеркал; выход из строя одного из дисков влечет потерю копии.
	\item \textit{Что такое LVM (Logical Volume Manager) и каковы его основные преимущества перед работой с обычными разделами?}\newline 
	LVM (Logical Volume Manager) -- это технология управления логическими томами, позволяющая объединить несколько физических дисков в один или несколько логических томов для более эффективного использования дискового пространства и упрощения управления ими.
	Динамическое изменение размеров дисков, разнесение томов по разным дискам, сохранение снапшотов и уровень абстракции между физическимм дисками и файловой системой.
	\item \textit{Что такое файловая система? Какие функции она предоставляет?}\newline 
	Файловая система -- способ организации, хранения, именования данных на носителях информации, правила доступа к файлам и каталогам. 
	Она предоставляет набор функций для работы с файлами (создание, копирование, перемещение, удаление, изменение размера и т.д.).
	\item \textit{Что такое "точка монтирования" (mount point) в Linux? Опишите шаги, необходимые для ручного монтирования раздела /dev/sdb1 с файловой системой ext4 в каталог /mnt/data.}\newline 
	Точка монтирования -- каталог в файловой системе, который становится корневым каталогом для подключаемой файловой системы. 
	Алгоритм создания: 
	создать точку монтирования (sudo mkdir /mnt/data)
	-> выполнить команду монтирования (sudo mount -t ext4 /dev/sda1 /mnt/data)
	-> проверить, что файловая система успешно примонтирована (df -h).
	\item \textit{Какой файл конфигурации в Linux отвечает за автоматическое монтирование файловых систем при загрузке?}\newline 
	Файл /etc/fstab.
\end{enumerate}
